% praktisk_proeve.tex -- indhold til praktisk prøve (bruges med \input{praktisk_proeve})
% Kræver i præamblen: \usepackage{longtable,booktabs,array,xcolor}
\providecolor{jagtgroen}{HTML}{2F4F3A}
\providecolor{jagtguld}{HTML}{B08D3C}
\newcommand{\fejlkode}[1]{\textcolor{jagtgroen}{\textbf{#1}}}

\section{Praktisk prøve -- geværhåndtering og skudafgivelse}

\subsection*{\textcolor{jagtgroen}{Fase 0: Generelle sikkerhedsmæssige forhold}}

\textbf{Korrekt håndtering}
\begin{itemize}\setlength{\itemsep}{2pt}
  \item Aftrækker og aftrækkerbøjle må ikke betjenes/berøres under prøven med undtagelse af under skudafgivelse (se Fase 6).
  \item Sikringsmekanismen må ikke betjenes/berøres under prøven med undtagelse af under skudafgivelse (se Fase 6).
  \item Aspiranten skal tydeligt orientere sig i den vandrette linie og sikre at området er fri for individer/genstande (f.eks. skure, biler, toiletter mm) inden geværet, i ladt tilstand, føres gennem den vandrette linie (begge veje).
  \item Geværløbene skal efterses for tilstopninger med himmel som baggrund hver gang aspiranten ikke kender tilstanden på geværet dvs. ved opstart, efter forcering af forhindring(er), hvis geværet overdrages til den prøvesagkyndige i forbindelse med iførelse af høreværn inden skydning.
  \item Tabes en eller flere klikpatroner ved ladning af geværet, tømmes geværet for eventuel isat klikpatron, patron(er) samles op fra jorden, og aspiranten starter forfra med at efterse geværløbene og lader først derefter geværet igen.
\end{itemize}

\begin{longtable}{@{}p{1.2cm}p{\dimexpr\textwidth-1.2cm-2\tabcolsep\relax}@{}}
\toprule
\textbf{Kode} & \textbf{Forkert håndtering} \\
\midrule
\endfirsthead
\toprule
\textbf{Kode} & \textbf{Forkert håndtering (fortsat)} \\
\midrule
\endhead
\bottomrule
\endlastfoot
\fejlkode{0A} & Fysisk kontakt eller tvivl om fysisk kontakt til aftrækker/aftrækkerbøjle. \\
\fejlkode{0B} & Fysisk kontakt eller tvivl om fysisk kontakt til sikringsmekanismen. \\
\fejlkode{0C} & Manglende tydelig orientering i den vandrette linie. \\
\fejlkode{0D} & Geværet lades (lukkes) uden patroner i kamrene. \\
\fejlkode{0E} & Geværløbene efterses ikke for tilstopninger. \\
\fejlkode{0F} & Geværløbene efterses for tilstopninger, men ikke med himmel som baggrund. \\
\end{longtable}
\bigskip

\subsection*{\textcolor{jagtgroen}{Fase 1: Geværet løftes fra stativet}}

\textbf{Korrekt håndtering}
\begin{itemize}\setlength{\itemsep}{2pt}
  \item Geværet løftes ved forskæftet med modsatte hånd af monteringsskulder-siden af kroppen.
  \item Geværet håndteres, så geværløbene peger lige op i luften og toplukkegrebet vender i retning af monteringsskulderen.
\end{itemize}

\noindent\colorbox{yellow!25}{%
  \begin{minipage}{\dimexpr\textwidth-2\fboxsep\relax}
    \textbf{Den prøvesagkyndiges involvering}\\
    Sagkyndige siger typisk: "Nu skal du vælge et gevær (fra stativet) og komme over til mig", dette kan forvirre aspiranten til at tro, at geværet skal bæres over til den sagkyndige i lukket tilstand (Kode 1B). Geværet skal håndteres efter instruktionen i Fase 2A el 2B. 
  \end{minipage}%
}
\bigskip

\begin{longtable}{@{}p{1.2cm}p{\dimexpr\textwidth-1.2cm-2\tabcolsep\relax}@{}}
\toprule
\textbf{Kode} & \textbf{Forkert håndtering} \\
\midrule
\endfirsthead
\toprule
\textbf{Kode} & \textbf{Forkert håndtering (fortsat)} \\
\midrule
\endhead
\bottomrule
\endlastfoot
\fejlkode{1A} & Geværet peges mod tagudhæng inden det knækkes. \\
\fejlkode{1B} & Geværet flyttes mere end 2 m fra stativ i lukket tilstand. \\
\end{longtable}
\bigskip
\bigskip

\subsection*{\textcolor{jagtgroen}{Fase 2A: Geværet knækkes (åbnes) med geværløbene i opadvendt retning}}

\textbf{Korrekt håndtering}
\begin{itemize}\setlength{\itemsep}{2pt}
  \item Geværet knækkes, så geværløbene vender opad og skæftet skydes til siden, så eventuelle patroner falder ud af geværet.
  \item Geværløbene efterses for tilstopninger med himmel som baggrund.
  \item Geværet kan herefter i knækket tilstand føres ned i en vilkårlig retning, dog væk fra den prøvesagkyndige (god skik).
  \item Geværet skal forblive knækket indtil det lades med klikpatroner, som udleveres af den prøvesagkyndige.
\end{itemize}

\begin{longtable}{@{}p{1.2cm}p{\dimexpr\textwidth-1.2cm-2\tabcolsep\relax}@{}}
\toprule
\textbf{Kode} & \textbf{Forkert håndtering} \\
\midrule
\endfirsthead
\toprule
\textbf{Kode} & \textbf{Forkert håndtering (fortsat)} \\
\midrule
\endhead
\bottomrule
\endlastfoot
\fejlkode{2A1} & Geværet knækkes med geværløbene pegende vandret eller tilnærmelsesvis vandret. \\
\fejlkode{0E} & Geværløbene efterses ikke for tilstopninger. \\
\fejlkode{0F} & Geværløbene efterses for tilstopninger, men ikke med himmel som baggrund. \\
\fejlkode{0D} & Geværet lades (lukkes) uden patroner i kamrene. \\
\end{longtable}
\bigskip

\subsection*{\textcolor{jagtgroen}{Fase 2B: Geværet knækkes (åbnes) med geværløbene i nedadvendt retning}}

\textbf{Korrekt håndtering}
\begin{itemize}\setlength{\itemsep}{2pt}
  \item Alternativ til Fase 2A.
  \item Efter tydelig orientering mod den vandrette linie føres geværet i lukket tilstand gennem den vandrette linie, så geværløbene peger skråt nedad.
  \item Geværet knækkes herefter og geværløbene efterses for tilstopninger med himmel som baggrund.
  \item Bemærk at denne metode øger risikoen for fejl da geværet skal føres unødvendigt ned gennem den vandrette linie sammenlignet med fremgangsmåden i Fase 2A.
  \item Der må ikke være fysisk kontakt eller tvivl om fysisk kontakt til sikringsmekanisme og aftrækker/aftrækkerbøjle.
\end{itemize}

\begin{longtable}{@{}p{1.2cm}p{\dimexpr\textwidth-1.2cm-2\tabcolsep\relax}@{}}
\toprule
\textbf{Kode} & \textbf{Forkert håndtering} \\
\midrule
\endfirsthead
\toprule
\textbf{Kode} & \textbf{Forkert håndtering (fortsat)} \\
\midrule
\endhead
\bottomrule
\endlastfoot
\fejlkode{0C} & Manglende tydelig orientering i den vandrette linie. \\
\fejlkode{0E} & Geværløbene efterses ikke for tilstopninger. \\
\fejlkode{0F} & Geværløbene efterses for tilstopninger, men ikke med himmel som baggrund. \\
\fejlkode{0D} & Geværet lades (lukkes) uden patroner i kamrene. \\
\end{longtable}
\bigskip

\subsection*{\textcolor{jagtgroen}{Fase 3: Ladning med klikpatroner}}

\textbf{Korrekt håndtering}
\begin{itemize}\setlength{\itemsep}{2pt}
  \item Den prøvesagkyndige udleverer klikpatroner og geværet lades (klikpatroner isættes) af aspiranten.
  \item Geværet må ikke flyttes til en ny retning efter at klikpatroner er isat selvom geværet stadig er knækket.
  \item Geværet lukkes i en tilpas fast bevægelse (ikke for hård / ikke for blød), hvor geværløbene peger væk fra aspirantens fødder, den vandrette linie samt den prøvesagkyndige.
  \item Skæftet lukkes op i geværløbene, som fastholdes i samme vinkel.
  \item Aspirantens fingre må ikke være ved eller i nærheden af sikringsmekanismen og aftrækker/aftrækkerbøjle.
\end{itemize}

\noindent\colorbox{yellow!25}{%
  \begin{minipage}{\dimexpr\textwidth-2\fboxsep\relax}
    \textbf{Den prøvesagkyndiges involvering}\\
    Sagkyndige kan placere sig på en uhensigsmæssig måde, set fra aspiranten, for at teste forståelsen af farligheden ved at lade geværet i knækket tilstand uden umiddelbart at kunne lukke geværet i en forsvarlig retning (Kode C1). 
  \end{minipage}%
}

\begin{longtable}{@{}p{1.2cm}p{\dimexpr\textwidth-1.2cm-2\tabcolsep\relax}@{}}
\toprule
\textbf{Kode} & \textbf{Forkert håndtering} \\
\midrule
\endfirsthead
\toprule
\textbf{Kode} & \textbf{Forkert håndtering (fortsat)} \\
\midrule
\endhead
\bottomrule
\endlastfoot
\fejlkode{C1} & Geværet flyttes til en ny retning efter at klikpatroner er isat i knækket tilstand. \\
\fejlkode{C2} & Geværet lukkes med geværløbene pegende i retning af den prøvesagkyndige. \\
\fejlkode{C3} & Geværet lukkes i eller nær den vandrette linie. \\
\fejlkode{C4} & Geværet lukkes i eller nær den lodrette linie (fødder). \\
\fejlkode{C5} & Geværet lukkes uforholdsmæssigt hårdt. \\
\fejlkode{C6} & Geværet lukkes uforholdsmæssigt blødt. \\
\end{longtable}
\bigskip

\subsection*{\textcolor{jagtgroen}{Fase 4: Geværet føres til aktiv jagtposition}}

\textbf{Korrekt håndtering}
\begin{itemize}\setlength{\itemsep}{2pt}
  \item Aspiranten orienterer sig tydeligt i den vandrette linie før geværet føres op i aktiv jagtposition.
  \item Vinklen på geværet skal være så andre personer trygt kan færdes rundt om aspiranten.
  \item Aspirantens fingre må ikke være ved eller i nærheden af sikringsmekanismen og aftrækker/aftrækkerbøjle.
\end{itemize}

\noindent\colorbox{yellow!25}{%
  \begin{minipage}{\dimexpr\textwidth-2\fboxsep\relax}
    \textbf{Den prøvesagkyndiges involvering}\\
    Sagkyndige kan placere sig på samme side som aspiranten peger geværløbene, hvis sagkyndige synes geværet bæres i en for flad vinkel (Kode 4A). 
  \end{minipage}%
}

\begin{longtable}{@{}p{1.2cm}p{\dimexpr\textwidth-1.2cm-2\tabcolsep\relax}@{}}
\toprule
\textbf{Kode} & \textbf{Forkert håndtering} \\
\midrule
\endfirsthead
\toprule
\textbf{Kode} & \textbf{Forkert håndtering (fortsat)} \\
\midrule
\endhead
\bottomrule
\endlastfoot
\fejlkode{0C} & Manglende tydelig orientering i den vandrette linie. \\
\fejlkode{4A} & Geværet bæres for lavt så geværløbene er nær den vandrette linie. \\
\fejlkode{0A} & Fysisk kontakt eller tvivl om fysisk kontakt til aftrækker/aftrækkerbøjle. \\
\fejlkode{0B} & Fysisk kontakt eller tvivl om fysisk kontakt til sikringsmekanismen. \\
\end{longtable}
\bigskip

\subsection*{\textcolor{jagtgroen}{Fase 5: Forcering af forhindringer: grøft og hegn}}

\textbf{Korrekt håndtering}
\begin{itemize}\setlength{\itemsep}{2pt}
  \item Den prøvesagkyndige giver instruks om at forhindring skal forceres.
  \item Aspiranten orienterer sig tydeligt i den vandrette linie hvorefter geværet føres ned så geværløbene peger væk fra fødder og vandret linie.
  \item Når geværet er i ro betjenes toplukkegrebet og geværet knækkes.
  \item Patroner tages ud og placeres i lettilgængelig lomme.
  \item Forhindringen forceres med knækket gevær.
  \item Geværløbene efterses for tilstopninger med himmel som baggrund.
  \item Aspiranten lader geværet med klikpatroner og lukker geværet i en fast bevægelse hvor geværløbene så vidt muligt er fikseret i samme retning.
  \item Aspiranten orienterer sig tydeligt i den vandrette linie og fører geværet op i aktiv jagtposition.
\end{itemize}

\begin{longtable}{@{}p{1.2cm}p{\dimexpr\textwidth-1.2cm-2\tabcolsep\relax}@{}}
\toprule
\textbf{Kode} & \textbf{Forkert håndtering} \\
\midrule
\endfirsthead
\toprule
\textbf{Kode} & \textbf{Forkert håndtering (fortsat)} \\
\midrule
\endhead
\bottomrule
\endlastfoot
\fejlkode{E1} & Geværet knækkes med geværløbene i opadvendt retning. \\
\fejlkode{E2} & Der orienteres ikke tydeligt før geværet føres gennem den vandrette linie. \\
\fejlkode{E3} & Geværet åbnes med geværløbene pegende i retning af den prøvesagkyndige. \\
\fejlkode{E4} & Geværet åbnes med geværløbene pegende nær den vandrette linie. \\
\fejlkode{E5} & Geværet åbnes med geværløbene pegende nær den lodrette linie. \\
\fejlkode{0E} & Geværløbene efterses ikke for tilstopninger. \\
\fejlkode{0F} & Geværløbene efterses for tilstopninger, men ikke med himmel som baggrund. \\
\fejlkode{C1} & Geværet flyttes til en ny retning efter at klikpatroner er isat i knækket tilstand. \\
\fejlkode{C2} & Geværet lukkes med geværløbene pegende i retning af den prøvesagkyndige. \\
\fejlkode{C3} & Geværet lukkes i eller nær den vandrette linie. \\
\fejlkode{C4} & Geværet lukkes i eller nær den lodrette linie (fødder). \\
\fejlkode{C5} & Geværet lukkes uforholdsmæssigt hårdt. \\
\fejlkode{C6} & Geværet lukkes uforholdsmæssigt blødt. \\
\end{longtable}
\bigskip

\subsection*{\textcolor{jagtgroen}{Fase 6: Skudafgivelse}}

\textbf{Korrekt håndtering}
\begin{itemize}\setlength{\itemsep}{2pt}
  \item Den prøvesagkyndige meddeler aspiranten, at der skal skydes til lerduer.
  \item Aspiranten orienterer sig tydeligt i den vandrette linie og fører geværet gennem den vandrette linie.
  \item Geværet knækkes og klikpatroner afleveres til den prøvesagkyndige.
  \item Geværet rækkes til den prøvesagkyndige, som holder geværet, imens aspiranten ifører sig høreværn.
  \item Aspiranten får geværet retur og gennemser geværløbene for tilstopninger med himmel som baggrund.
  \item Aspiranten lader geværet med 2 haglpatroner.
  \item Aspiranten fører geværet op til aktiv jagtposition. Geværløbene skal pege opad i en vinkel så en person kan bevæge sig foran aspiranten.
  \item Aspiranten giver verbalt meldingen: ``Klar''.
  \item Aspiranten monterer geværet i skulderen og afsikrer geværet (sikringen skubbes frem).
  \item Aspiranten skal tydeligt have kontakt mellem geværskæfte og kind ved skudafgivelse.
  \item Den prøvesagkyndige skal aktivere lerduekasteren ikke længere end 5 sek efter aspirantens klarmelding.
  \item Aspiranten skyder ét skud i retning af lerduen (bagduen).
  \item Aspiranten sikrer geværet (sikringen skubbes tilbage).
  \item Aspiranten fører geværet ned til aktiv jagtposition.
  \item Aspiranten fører geværet ned gennem den vandrette linie.
  \item Aspiranten knækker geværet og tager afskudt patron ud.
  \item Aspiranten lader det tomme kammer med en ny patron og gentager proceduren ved skudafgivelse.
\end{itemize}

\begin{longtable}{@{}p{1.2cm}p{\dimexpr\textwidth-1.2cm-2\tabcolsep\relax}@{}}
\toprule
\textbf{Kode} & \textbf{Forkert håndtering} \\
\midrule
\endfirsthead
\toprule
\textbf{Kode} & \textbf{Forkert håndtering (fortsat)} \\
\midrule
\endhead
\bottomrule
\endlastfoot
\fejlkode{F1} & Geværet lukkes efter at klikpatroner er afleveret. \\
\fejlkode{F2} & Geværløbene efterses ikke for tilstopninger efter at høreværn er iført. \\
\fejlkode{F3} & Vandret geværføring (montering) inden klarmelding. \\
\fejlkode{F4} & Sikring betjenes inden klarmelding. \\
\fejlkode{F5} & Aftrækker betjenes inden klarmelding. \\
\fejlkode{F6} & Geværet afsikres ikke før skudafgivelse. \\
\fejlkode{F7} & Manglende kontakt mellem kind og skæfte ved skudafgivelse. \\
\fejlkode{F8} & Toplukkegreb betjenes inden geværet er nedadpegende. \\
\end{longtable}
\bigskip

\subsection*{\textcolor{jagtgroen}{Fase 7: Aflevering af gevær}}

\textbf{Korrekt håndtering}
\begin{itemize}\setlength{\itemsep}{2pt}
  \item Efter det 4. skud til lerdue tømmer aspiranten geværet for patroner.
  \item Den afskudte patron smides ud som de forrige 3 patroner.
  \item Den ikke afskudte patron placeres i kassen/kurven (samme sted som den blev taget fra).
  \item Geværet rækkes til den prøvesagkyndige med skæftet forrest, så den prøvesagkyndige kan se, at geværet er tomt for ammunition.
\end{itemize}

\noindent\colorbox{yellow!25}{%
  \begin{minipage}{\dimexpr\textwidth-2\fboxsep\relax}
    \textbf{Den prøvesagkyndiges involvering}\\
    Umiddelbart efter at aspiranten har affyret sidste skud til lerdue, siger sagkyndige typisk: "Nu kan du aflevere geværet til mig". Her tester sagkyndige om aspiranten forstår, at geværet skal tømmes for afskudt og intakt ammunition inden det overdrages. 
  \end{minipage}%
  }

\begin{longtable}{@{}p{1.2cm}p{\dimexpr\textwidth-1.2cm-2\tabcolsep\relax}@{}}
\toprule
\textbf{Kode} & \textbf{Forkert håndtering} \\
\midrule
\endfirsthead
\toprule
\textbf{Kode} & \textbf{Forkert håndtering (fortsat)} \\
\midrule
\endhead
\bottomrule
\endlastfoot
\fejlkode{7A} & Geværet afleveres med en eller flere patroner i kammeret/kamrene. \\
\end{longtable}
\bigskip